\chapter*{Abstract}
\addcontentsline{toc}{chapter}{Abstract}
\markboth{Abstract}{Abstract}

This book develops a unified mathematical framework for periodic phenomena, based on the principle ``symmetry $\to$ invariant $\to$ closed orbit $\to$ period''.

The starting point is that periodic phenomena in different disciplines---celestial motion, electrical oscillation, biological rhythm, economic fluctuation and algorithmic iteration---have long been described in mutually untranslatable languages. The book argues that they can all be captured by a single mathematical object: a \textbf{family of epicycles}. A system of period $T$ is described by a set of amplitude, angular-frequency and phase triples $(A_k,\omega_k,\phi_k)$, whilst the relations between periods are captured by \textbf{periodic convolution}. Periodicity is thereby reduced to a question about group actions and group representations: the three basic transformations---translation, scaling and rotation---generate every elementary periodic structure, and their completeness can be proved rigorously within the framework of one-parameter families of conformal transformations of the complex plane.

The core methodology consists of four pillars: group representation theory, periodic convolution, window functions and refraction-index diagnostics, and complex exponential activation functions (CEAF) with kernel-smoothed derivative--integral descent. On this basis the book gives a standard five-step interface for attaching an arbitrary target to the framework, and applies it domain by domain in more than twenty fields, from astronomy, physics, biology and psychology to economics, cryptography and computer science. Throughout, three things are kept strictly apart: theorems that have been proved, conceptions advanced as research programmes, and problems that remain open.

\vspace{1em}
\noindent\textbf{Keywords:} periodic mathematics; epicycle; harmonic analysis; conservation law; symmetry; phase; amplitude; frequency; periodic convolution; group theory; Fourier transform; Mellin transform; wavelet transform; I Ching; Leibniz; calendar; chip; signal processing

\newpage
